\documentclass[conference]{IEEEtran}
\IEEEoverridecommandlockouts

\usepackage{cite}
\usepackage{amsmath,amssymb,amsfonts}
\usepackage{algorithmic}
\usepackage{graphicx}
\usepackage{textcomp}
\usepackage{xcolor}
\usepackage{booktabs}
\usepackage{array}
\usepackage{multirow}

\def\BibTeX{{\rm B\kern-.05em{\sc i\kern-.025em b}\kern-.08em
    T\kern-.1667em\lower.7ex\hbox{E}\kern-.125emX}}

\begin{document}

\title{Task-Driven Object Detection with Semantic Distillation and Dynamic Gating
Using an Accelerator on FPGA for Embedded Applications}

\author{%
\IEEEauthorblockN{Ackshaya Keerthi G}
\IEEEauthorblockA{\textit{Dept. of Computer Science \& Engineering}\\
\textit{RV College of Engineering}\\
Bengaluru, India\\
1rv23cs013@rvce.edu.in}
\and
\IEEEauthorblockN{Bhoomika Sundar}
\IEEEauthorblockA{\textit{Dept. of Computer Science \& Engineering}\\
\textit{RV College of Engineering}\\
Bengaluru, India\\
1rv23cs066@rvce.edu.in}
\and
\IEEEauthorblockN{Tandle Suhani}
\IEEEauthorblockA{\textit{Dept. of Electronics \& Communication}\\
\textit{RV College of Engineering}\\
Bengaluru, India\\
1rv22ec184@rvce.edu.in}
\and
\IEEEauthorblockN{Vaibhavi D}
\IEEEauthorblockA{\textit{Dept. of Electronics \& Communication}\\
\textit{RV College of Engineering}\\
Bengaluru, India\\
1rv22ec177@rvce.edu.in}
\and
\IEEEauthorblockN{Dr. Roopa T S}
\IEEEauthorblockA{\textit{Dept. of Mechanical Engineering}\\
\textit{RV College of Engineering}\\
Bengaluru, India\\
roopaTS@rvce.edu.in}
}

\maketitle

% ---------------------------------------------------------------
\begin{abstract}
Current object detection systems for edge devices are constrained by the computational
burden of cross-modal attention mechanisms, which carry a quadratic complexity of
$O(N^{2})$. These architectures produce large similarity matrices and rely on Softmax
functions that are incompatible with the fixed-point arithmetic of embedded processors
such as the VEGA RISC-V processor, creating critical memory bottlenecks. This paper
proposes a task-driven pipeline that replaces heavy cross-modal attention with a
lightweight dynamic gating mechanism operating at the channel level, reducing fusion
complexity to $O(N)$. A YOLOv8 visual backbone is distilled from a frozen CLIP teacher
model, transferring functional affordance reasoning into a compact 6\,MB student model
via Mean Squared Error minimization. A Task-Mapping Multi-Layer Perceptron (MLP)
generates Hadamard-product gating coefficients that selectively activate or suppress
feature channels based on semantic intent. Deployment on the Genesys-2 FPGA using
INT8 quantization and bit-shift arithmetic yields a 96.29\% average channel sparsity,
a 40.93\% improvement in semantic alignment over the YOLOv8 baseline, and a
77.04\% estimated dynamic power saving---all while maintaining 100\% detection
consistency with the dense baseline across 5,000 test images.
\end{abstract}

\begin{IEEEkeywords}
FPGA, object detection, knowledge distillation, dynamic gating, semantic affordance,
YOLOv8, CLIP, INT8 quantization, edge computing, VEGA processor
\end{IEEEkeywords}

% ---------------------------------------------------------------
\section{Introduction}

The proliferation of embedded vision systems in robotics, assistive devices, and
autonomous agents demands a shift from general-purpose object detection toward
\emph{task-aware} perception. Rather than enumerating every object in a scene, a
task-driven system must interpret a natural-language intent (e.g., ``find something
to sit on'') and prioritize only the relevant visual affordances. This is especially
critical for resource-constrained hardware where energy and memory budgets are
strictly bounded.

Standard cross-modal attention, as used by models such as GLIP~\cite{li2022grounded}
and YOLO-World~\cite{cheng2024yoloworld}, achieves strong semantic grounding but
incurs $O(N^{2})$ complexity and floating-point operations (FP32) that exceed the
capabilities of fixed-point embedded processors. Specifically, the VEGA processor
on the Genesys-2 FPGA board lacks hardware support for the Softmax exponential
required by standard attention, causing memory overflows in on-chip Block RAM (BRAM).

To bridge this gap, we propose a dual-encoder pipeline that (i) distils functional
affordance knowledge from a CLIP~\cite{radford2021clip} teacher into a lightweight
YOLOv8 student, and (ii) replaces attention with an $O(N)$ \emph{dynamic gating}
module implemented as element-wise Hadamard products. The complete system is
quantized to INT8 and deployed on the CDAC VEGA RISC-V processor with bit-shift
arithmetic for real-time inference.

The main contributions of this paper are:
\begin{itemize}
    \item A task-conditional dynamic sparsity formulation that reduces cross-modal
          fusion complexity from $O(N^{2})$ to $O(N)$.
    \item A CLIP-to-YOLOv8 semantic distillation pipeline that achieves a 40.93\%
          reduction in affordance localization error.
    \item An INT8 hardware deployment achieving 96.29\% average channel sparsity
          and 77.04\% estimated dynamic power saving on the VEGA processor, with
          zero detection consistency loss.
    \item A SystemVerilog verification framework confirming all CNN accelerator
          test cases pass on the Genesys-2 FPGA.
\end{itemize}

% ---------------------------------------------------------------
\section{Related Work}

\subsection{Vision-Language Models for Detection}
Radford \emph{et al.}~\cite{radford2021clip} introduced CLIP, enabling zero-shot
object identification via contrastive language-image pre-training. Li \emph{et
al.}~\cite{li2022grounded} extended this to grounded detection with GLIP, framing
detection as a phrase-grounding problem. Both models, however, rely on dense
transformer attention whose quadratic complexity and FP32 requirements preclude
direct FPGA deployment.

\subsection{Efficient Distillation and Compact Models}
Vasu \emph{et al.}~\cite{vasu2023mobileclip} demonstrated with MobileCLIP that
teacher-student distillation can compress large vision-language models to
mobile-scale footprints while preserving much of the teacher's semantic
reasoning. This motivates our use of CLIP as the affordance teacher for YOLOv8.

\subsection{Dynamic and Sparse Computation}
Yang~\cite{yang2022rosetta} utilized task-aware gates in ROSETTA to activate
task-specific channels; Rao \emph{et al.}~\cite{rao2021dynamicvit} developed
DynamicViT to prune irrelevant image tokens at inference time. Neither work
addresses INT8 FPGA deployment. Sawatzky \emph{et al.}~\cite{sawatzky2019affordance}
applied Gated Graph Neural Networks to region prioritization, but message-passing
overhead is prohibitive for real-time edge silicon. Our approach unifies channel-level
gating with hardware-optimized quantization to fill this gap.

% ---------------------------------------------------------------
\section{Methodology}

\subsection{System Overview}

\begin{figure}[htbp]
\centerline{\fbox{\rule{0pt}{3.5cm}\rule{8cm}{0pt}}}
\caption{High-level dual-encoder pipeline. Natural language task input is
encoded into a 512-dimensional embedding that drives the dynamic gating module
acting on the YOLOv8 P5 feature map before the INT8 quantization layer.}
\label{fig:pipeline}
\end{figure}

The proposed pipeline (Fig.~\ref{fig:pipeline}) ingests a $640\times640$ RGB
image and a natural-language task string simultaneously. A frozen CLIP text encoder
maps the task to a 512-dimensional embedding $\mathbf{z}_{\text{task}}$. A YOLOv8n
backbone extracts the P5 feature map $\mathbf{F}\in\mathbb{R}^{256\times20\times20}$.
A lightweight Task-Mapping MLP converts $\mathbf{z}_{\text{task}}$ into a gating
vector $\mathbf{g}\in\mathbb{R}^{256}$, which is applied to $\mathbf{F}$ via a
Hadamard product before the detection head.

\subsection{Task-Conditional Dynamic Gating}

The core operation of the gating layer at layer $l$ is:
\begin{equation}
    Y_l = \sigma\!\left(\mathcal{G}(Z_{\text{task}};\,W_g)\right) \odot f(X_{l-1};\,W_l),
    \label{eq:gating}
\end{equation}
where $\mathcal{G}$ is the gating MLP, $\sigma$ is the sigmoid activation, $\odot$
denotes the element-wise Hadamard product, $X_{l-1}$ is the visual input from the
preceding layer, and $W_g$, $W_l$ are learned weight matrices. The MLP transformation
producing the gating coefficients is:
\begin{equation}
    \mathbf{g} = \sigma\!\left(W_2 \cdot \phi(W_1 \cdot \mathbf{v} + \mathbf{b}_1) + \mathbf{b}_2\right),
    \label{eq:mlp}
\end{equation}
where $\phi$ is ReLU and $\mathbf{v}\in\mathbb{R}^{64}$ is the task embedding
vector. Sigmoid constrains each coefficient $g_i\in(0,1)$, with values near zero
indicating that the corresponding convolutional channel is semantically redundant
and can be skipped by the hardware.

The modulated feature map applied to channel $c$ at spatial position $(h,w)$ is:
\begin{equation}
    F'_{c,h,w} = F_{c,h,w} \times g_c.
    \label{eq:hadamard}
\end{equation}

\subsection{Sparsity-Inducing Training Objective}

To encourage hardware-friendly sparsity, the total training loss incorporates an
$L_1$ regulariser on gating activations:
\begin{equation}
    \mathcal{L}_{\text{total}} = \mathcal{L}_{\text{detection}} + \lambda \sum_{i}\left|\sigma\!\left(\mathcal{G}(Z_{\text{task}})\right)_i\right|.
    \label{eq:loss}
\end{equation}
This penalises non-zero gates, incentivising the model to deactivate as many
channels as possible without degrading detection accuracy.

\subsection{Semantic Distillation from CLIP}

The distillation loss aligns the student's gated feature activations with the
CLIP teacher's attention maps:
\begin{equation}
    \mathcal{L}_{\text{distill}} = \frac{1}{N}\sum_{i=1}^{N}
        \left\|\Phi_T(I_i, T_{\text{text}}) - \Phi_S(I_i, \mathbf{g})\right\|^2,
    \label{eq:distill}
\end{equation}
where $\Phi_T$ and $\Phi_S$ are the teacher and student feature extraction
functions for image $I_i$ and text query $T_{\text{text}}$. Minimising this
Euclidean distance transfers the CLIP teacher's 2\,GB semantic intelligence into
the 6\,MB student without requiring CLIP at inference time.

\subsection{INT8 Quantization and Hardware Deployment}

All floating-point weights are quantized to 8-bit integers using:
\begin{equation}
    Q(x) = \text{clamp}\!\left(\left\lfloor\frac{x}{S}\right\rfloor + Z,\,-128,\,127\right),
    \label{eq:quant}
\end{equation}
where $S$ is the per-tensor scale factor and $Z$ is the zero-point offset. The
gating coefficients are further optimised using bit-shift arithmetic, replacing
FP32 multiplications with fast integer shifts on the VEGA RISC-V ISA. The
convolution engine checks each gate $g_c$ against a threshold $\tau=0.1$ before
issuing a Multiply-Accumulate (MAC) instruction; if $g_c<\tau$, the entire
feature plane is bypassed, reducing dynamic switching activity.

\subsection{Dataset and Task Mapping}

Training and evaluation use the COCO-2017 dataset. A 14-task affordance vocabulary
(e.g., \textit{sitting}, \textit{grasping}, \textit{pouring}, \textit{cutting}) is
defined, and each task is mapped to relevant COCO object categories. Input images
are resized to $640\times640$ pixels and normalised to $[0,1]$. During distillation,
CLIP generates reference attention maps stored as ground-truth tensors for the
student to minimise~\eqref{eq:distill}.

% ---------------------------------------------------------------
\section{Verification of the CNN Accelerator}

A layered SystemVerilog testbench validates the CNN accelerator hardware before
FPGA deployment. The verification architecture follows a modular
Driver~$\to$~DUT~$\to$~Monitor~$\to$~Scoreboard pipeline, exercising the design
under three test cases:

\begin{enumerate}
    \item \textbf{Normal Inference}: Reset $\to$ Start $\to$ CNN processing $\to$ Done;
          verifies FSM transitions, Conv1, pooling, Conv2, and fully connected outputs.
    \item \textbf{Mid-Run Reset Recovery}: Start $\to$ wait 30 cycles $\to$ assert reset
          $\to$ restart; verifies counter and BRAM recovery.
    \item \textbf{Back-to-Back Inference}: Run\,\#1 $\to$ Reset $\to$ Run\,\#2; confirms
          accumulator and memory clearing between consecutive inferences.
\end{enumerate}

The DUT (\texttt{top\_cnn}) contains an input BRAM, Conv1 engine with ReLU, a
ping-pong BRAM and max-pool unit, a Conv2 layer, a fully connected layer, and an
FSM controller. A 100\,MHz clock is generated by \texttt{always \#5 clk = \textasciitilde clk}.
The scoreboard validates the fully connected outputs against the golden reference
\texttt{fc = [401, 2, 403, 4]}, yielding a final result of $\text{PASS}=4$,
$\text{FAIL}=0$---all tests passed.

% ---------------------------------------------------------------
\section{Results and Discussion}

\subsection{Semantic Alignment (Simulation)}

Semantic alignment is measured by the Mean Squared Error (MSE) between the
CLIP teacher heatmap $Y_{i,j}$ and the student activation map $\hat{Y}_{i,j}$:
\begin{equation}
    \text{MSE} = \frac{1}{M \times N}\sum_{i=1}^{M}\sum_{j=1}^{N}
        \left(Y_{i,j} - \hat{Y}_{i,j}\right)^2.
    \label{eq:mse}
\end{equation}

Table~\ref{tab:alignment} reports results over 5,000 COCO test images.

\begin{table}[htbp]
\caption{Semantic Alignment Performance over 5,000 Test Samples}
\label{tab:alignment}
\begin{center}
\begin{tabular}{lcc}
\toprule
\textbf{Model} & \textbf{MSE} & \textbf{Improvement} \\
\midrule
Standard YOLOv8 Baseline & 0.1343 & --- \\
Task-Aware Student Model & 0.0710 & 40.93\% \\
\bottomrule
\end{tabular}
\end{center}
\end{table}

The task-aware student achieves a 40.93\% reduction in MSE, confirming that the
CLIP-guided distillation successfully transfers functional affordance reasoning.
The error distribution for the student model (Fig.~\ref{fig:mse_dist}) is
left-shifted and sharper than the baseline, indicating consistent localization
across diverse scenes rather than isolated improvements.

\begin{figure}[htbp]
\centerline{\fbox{\rule{0pt}{3cm}\rule{8cm}{0pt}}}
\caption{Probability density distribution of MSE scores for the standard YOLOv8
baseline (pink) and the task-aware student model (green). The student exhibits a
sharper, left-skewed distribution, indicating more consistent semantic alignment.}
\label{fig:mse_dist}
\end{figure}

\subsection{Qualitative Affordance Localization}

Side-by-side heatmap comparisons for the \emph{sitting} and \emph{grasping} tasks
demonstrate that the student model (Fig.~\ref{fig:heatmaps}) faithfully replicates
the CLIP teacher's focus. For \emph{sitting}, the student highlights horizontal
support surfaces (chair seats, sofas), ignoring walls and floors. For
\emph{grasping}, the model concentrates activations on graspable edges and handles
rather than the entire object, enabling fine-grained interaction guidance.

\begin{figure}[htbp]
\centerline{\fbox{\rule{0pt}{3.5cm}\rule{8cm}{0pt}}}
\caption{Qualitative affordance localization. Left: gated detections. Centre: CLIP
teacher focus (ground truth). Right: student model focus. Shown for the
\emph{sitting} task (top row) and \emph{grasping} task (bottom row).}
\label{fig:heatmaps}
\end{figure}

\subsection{Hardware Gating Sparsity}

Channel sparsity is defined as:
\begin{equation}
    S_c = \left(1 - \frac{\sum_{i=1}^{C}\mathbb{1}(g_i > \tau)}{C}\right)\times 100,
    \label{eq:sparsity}
\end{equation}
where $C=256$, $\tau=0.1$. Table~\ref{tab:sparsity} summarises hardware
efficiency metrics averaged across all 14 tasks.

\begin{table}[htbp]
\caption{Hardware Sparsity and Theoretical Efficiency Metrics}
\label{tab:sparsity}
\begin{center}
\begin{tabular}{lc}
\toprule
\textbf{Metric} & \textbf{Value} \\
\midrule
Average Channel Sparsity $S_c$ & 96.29\% \\
Active Computational Channels & 9.47 / 256 \\
Theoretical Computational Reduction & 96.30\% \\
Estimated Dynamic Power Saving (VEGA) & 77.04\% \\
\bottomrule
\end{tabular}
\end{center}
\end{table}

Fig.~\ref{fig:sparsity_bar} shows per-task sparsity, consistently exceeding 80\%
for all 14 tasks and averaging 96.29\%. The gating correlation analysis
(Pearson $\rho=-0.1298$) reveals that the Task Mapper learns a deterministic,
image-invariant mask per task. This is advantageous for FPGA deployment: gating
configurations can be pre-loaded into BRAM at task-selection time, eliminating
real-time MLP overhead during inference.

\begin{figure}[htbp]
\centerline{\fbox{\rule{0pt}{3cm}\rule{8cm}{0pt}}}
\caption{Percentage of YOLOv8 P5 channels deactivated per semantic task.
The dashed red line marks the 96.3\% mean across all 14 tasks.}
\label{fig:sparsity_bar}
\end{figure}

\subsection{Inter-Task Semantic Specialization}

The Jaccard Similarity Index between active channel masks for any two tasks:
\begin{equation}
    J(T_1, T_2) = \frac{|\mathbf{G}_{T_1} \cap \mathbf{G}_{T_2}|}{|\mathbf{G}_{T_1} \cup \mathbf{G}_{T_2}|},
    \label{eq:jaccard}
\end{equation}
yields a mean inter-task overlap of just 1.85\%. Semantically proximate tasks
(\textit{carrying}--\textit{balancing}) peak at 0.82, while structurally distinct
tasks approach zero overlap. This confirms high architectural specialization: the
hardware activates a nearly unique set of channels per task.

\subsection{Functional Integrity and Detection Consistency}

\begin{table}[htbp]
\caption{Detection Consistency Across 5,000 Images}
\label{tab:consistency}
\begin{center}
\begin{tabular}{lc}
\toprule
\textbf{Metric} & \textbf{Value} \\
\midrule
Detection Consistency Rate & 100.00\% \\
Mean Object Recall Parity & 1.00 \\
Person Detection Recall Stability & 100.00\% \\
\bottomrule
\end{tabular}
\end{center}
\end{table}

Despite 96.29\% channel suppression, the task-aware model maintains 100\%
detection consistency with the dense YOLOv8 baseline (Table~\ref{tab:consistency}).
The suppressed channels in the P5 layer are semantically redundant for object
localization, confirming that gating operates as a semantic filter without
degrading foundational detection capability.

\subsection{Inference Latency and Software Overhead}

\begin{table}[htbp]
\caption{Inference Latency Benchmarking (NVIDIA T4 GPU, 100 Trials)}
\label{tab:latency}
\begin{center}
\begin{tabular}{lcc}
\toprule
\textbf{Metric} & \textbf{YOLOv8n} & \textbf{Task-Aware} \\
\midrule
Mean Inference Latency (ms) & 8.17 & 9.12 \\
Software Overhead $\Omega$ & --- & 11.65\% \\
\bottomrule
\end{tabular}
\end{center}
\end{table}

The 11.65\% software overhead (Table~\ref{tab:latency}) arises from sequential
Task Mapper and Gating MLP execution on a general-purpose GPU. On the VEGA
RISC-V processor, the custom ISA leverages sparsity to bypass MAC operations,
translating the 96.30\% mathematical reduction in the P5 layer into substantial
hardware-level power savings that far outweigh the software overhead.

% ---------------------------------------------------------------
\section{Conclusion}

We presented a task-driven object detection system that achieves $O(N)$ cross-modal
fusion complexity through dynamic channel gating, enabling deployment on FPGA-based
edge hardware that is incompatible with standard attention mechanisms. By distilling
CLIP's semantic affordance knowledge into a compact YOLOv8 student and replacing
similarity matrices with Hadamard-product masks, the proposed pipeline realizes:
(i) a 40.93\% semantic alignment improvement over the dense baseline;
(ii) 96.29\% average channel sparsity with an estimated 77.04\% dynamic power saving
on the VEGA processor;
(iii) 100\% detection consistency with the dense model across 5,000 test images; and
(iv) full hardware verification of the CNN accelerator via SystemVerilog, with all
four test scenarios passing.

Future work will extend the task vocabulary beyond the current 14 predefined prompts
to open-vocabulary queries processed in real time, investigate hybrid INT4/INT8 
quantization, and explore ASIC tape-out to further reduce power consumption for
wearable assistive devices.

% ---------------------------------------------------------------
\section*{Acknowledgment}

The authors thank Dr.\ Roopa T.S.\ for guidance throughout this work, Dr.\ Kariyappa
for panel review, and Prof.\ Narashimaraja P.\ for the organized \LaTeX{} template.
This work was carried out at RV College of Engineering, Bengaluru, under the
interdisciplinary project programme (XX367P) for the academic year 2025--26.

% ---------------------------------------------------------------
\begin{thebibliography}{00}

\bibitem{radford2021clip}
A.\ Radford, J.\ W.\ Kim, C.\ Hallacy, A.\ Ramesh, G.\ Goh, S.\ Agarwal,
G.\ Sastry, A.\ Askell, P.\ Mishkin, J.\ Clark, G.\ Krueger, and I.\ Sutskever,
``Learning transferable visual models from natural language supervision,''
in \emph{Proc.\ Int.\ Conf.\ Mach.\ Learn.\ (ICML)}, 2021, pp.\ 8748--8763.

\bibitem{li2022grounded}
L.\ Li, P.\ Zhang, H.\ Zhang, J.\ Yang, C.\ Li, Y.\ Zhong, L.\ Wang, L.\ Yuan,
L.\ Zhang, J.\ N.\ Hwang, K.\ Chang, and J.\ Gao,
``Grounded language-image pre-training,''
in \emph{Proc.\ IEEE/CVF Conf.\ Comput.\ Vis.\ Pattern Recognit.\ (CVPR)},
2022, pp.\ 10965--10975.

\bibitem{cheng2024yoloworld}
T.\ Cheng, L.\ Song, Y.\ Ge, W.\ Liu, X.\ Wang, and Y.\ Shan,
``YOLO-World: Real-time open-vocabulary object detection,''
in \emph{Proc.\ IEEE/CVF Conf.\ Comput.\ Vis.\ Pattern Recognit.\ (CVPR)},
2024, pp.\ 16901--16911.

\bibitem{vasu2023mobileclip}
P.\ K.\ A.\ Vasu, H.\ Fathi, J.\ Gao, O.\ Tuzel, and A.\ Ranjan,
``MobileCLIP: Fast image-text models through multi-modal reinforced training,''
in \emph{Proc.\ IEEE/CVF Conf.\ Comput.\ Vis.\ Pattern Recognit.\ (CVPR)},
2024, pp.\ 15963--15972.

\bibitem{yang2022rosetta}
Z.\ Yang, Y.\ Zhu, Y.\ Wu, M.\ Bain, F.\ Zhan, J.\ Li, and B.\ Shi,
``ROSETTA: Self-supervised speech and text joint learning for spoken language
understanding,'' \emph{arXiv preprint arXiv:2204.01581}, 2022.

\bibitem{rao2021dynamicvit}
Y.\ Rao, W.\ Zhao, B.\ Liu, J.\ Lu, J.\ Zhou, and C.-J.\ Hsieh,
``DynamicViT: Efficient vision transformers with dynamic token sparsification,''
in \emph{Adv.\ Neural Inf.\ Process.\ Syst.\ (NeurIPS)}, vol.\ 34, 2021,
pp.\ 13937--13949.

\bibitem{sawatzky2019affordance}
J.\ Sawatzky, A.\ Srikantha, and J.\ Gall,
``Affordance detection in RGB-D images,''
in \emph{Proc.\ IEEE Conf.\ Comput.\ Vis.\ Pattern Recognit.\ Workshops
(CVPRW)}, 2019, pp.\ 0--0.

\bibitem{jocher2023yolov8}
G.\ Jocher, A.\ Chaurasia, and J.\ Qiu,
``Ultralytics YOLOv8,'' 2023. [Online].
Available: \url{https://github.com/ultralytics/ultralytics}

\bibitem{lin2014coco}
T.-Y.\ Lin, M.\ Maire, S.\ Belongie, L.\ Bourdev, R.\ Girshick, J.\ Hays,
P.\ Perona, D.\ Ramanan, C.\ L.\ Zitnick, and P.\ Dollar,
``Microsoft COCO: Common objects in context,''
in \emph{Proc.\ Eur.\ Conf.\ Comput.\ Vis.\ (ECCV)}, 2014, pp.\ 740--755.

\end{thebibliography}

\end{document}
